\section{Research Question and Motivation}
A regression test must encode the reported behavior, execute in the project's native suite, and distinguish buggy from corrected code. An agent can produce a plausible assertion that passes prematurely, fails for an unrelated reason, or becomes invalid after a legitimate API change. SWT-Bench evaluates repository-level generated tests against reference fixes and supplies complementary behavioral and coverage evidence~\cite{swtpaper,swtrepo}.

We use DeepSeek-Flash for test generation in both workflows. We ask: \emph{Does evidence-grounded issue-specification extraction and validation improve DeepSeek-Flash's issue reproduction, as measured by official SWT-Bench instance-level success, compared with direct report-to-test generation?} TRACE (Test Reproduction through Analysis, Checks, and Evidence) is motivated by the hypothesis that explicitly identifying a target, trigger, observed behavior, expected behavior, and suggested assertion will reduce unsupported assumptions. We measure reproduction success, not the number of unsupported assumptions, so the proposed mechanism is not directly tested. Validation is intended to filter the issue specification before generation; it does not independently prove test correctness. Our contribution is a bounded preprocessing workflow, preserved paired artifacts, and an empirical examination of when structured issue specifications help or fail.

\section{Method and Experimental Design}\label{sec:method}
\subsection{Two Workflows}
In this paper, ``DeepSeek'' means DeepSeek-Flash; recorded provider identifiers appear in Appendix~\ref{app:disclosure}.

\textbf{Baseline.} Our Python script, \code{controller.py}, reads the saved issue description from \code{report.txt}. It passes that text unchanged to mini-SWE-agent 2.4.6~\cite{mini} by calling \code{agent.run(task)}. The controller does not summarize or analyze the issue first. mini-SWE-agent combines the issue with the shared instructions for writing regression tests and sends these messages to DeepSeek. DeepSeek proposes shell commands to inspect source code, read existing tests, write a test, and run it. mini-SWE-agent executes those commands in a Docker container containing the buggy repository, then sends their output and exit status back to DeepSeek. This exchange repeats as the model examines results and adjusts its test. DeepSeek selects the actions; mini-SWE-agent carries them out and maintains the conversation. The controller saves the final submitted patch and the conversation record.

\textbf{TRACE.} Before starting that same agent process, the controller makes a separate request to DeepSeek to analyze only the original issue. It asks for five items: the target function, triggering input, observed behavior, expected behavior, and a suggested assertion. Each known item must include an exact supporting quote from the issue; an unknown item and its evidence must both be marked null. Python checks that the response has the required structure, that its quotes occur in the report, and that at least one item contains usable information. A second, separate DeepSeek request compares the extracted information with the original issue description. It checks whether the issue supports each item, whether any item contradicts the issue, and whether important information was left out. These checks assess the issue specification, not the generated test.

If either check rejects the issue specification, the controller sends the errors back with the original issue and requests another extraction. It permits up to three extraction attempts; each attempt that passes the Python checks is followed by a review, giving at most six preprocessing requests. When an issue specification is accepted, the controller appends it to the original issue and passes both to mini-SWE-agent. If all attempts are rejected, it passes only the original issue and records a fallback. Missing credentials and other operational errors stop the run rather than triggering this fallback. The agent then uses the same DeepSeek-driven inspection, test-writing, execution, and submission process as the baseline. The specification-extraction and review requests are separate from the agent's generation conversation; only the accepted issue specification is handed from one to the other.

The course guidelines require a same-model workflow comparison and official unit-test evaluation~\cite{guidelines}. The preserved run records support this procedure~\cite{artifacts}. Both methods use the same native-test instructions, buggy environment, and generation limits: 100 agent steps and a nominal \$0.50 agent budget, excluding TRACE preprocessing. A final model request can take spending above that nominal limit. The agent's submitted patch is exported unchanged. All 52 selected runs submitted patches, so no empty predictions occurred in the final evaluated set; the exporter retains non-submissions as empty predictions when they occur. Neither method receives the reference fix or evaluator-only results during generation or repair. Running a test on buggy code supplies feedback, but does not establish that it passes after the fix. An automatic classifier that distinguishes issue-related failures from syntax, import, or setup errors is future work, not an implemented controller step.

The original agent configuration asks for a standalone script that reproduces the bug. For this experiment, both methods instead receive the same instructions to add regression tests to the repository's existing test suite. Both also enable Bash's \code{pipefail} option, which preserves a failing test runner's exit status when its output is piped through another command. These shared changes make the outputs suitable for unit-test evaluation and keep execution failures visible. The original configuration file is preserved; the controller applies these changes in memory for each run. The comparison therefore measures the addition of TRACE's issue specification to this shared unit-test workflow.

\subsection{Assignment, Controls, and Evaluation}
The fixed assignment comprises 26 issues, 13 Easy and 13 Hard, across 12 Python repositories. These are historical assignment labels, not our measured difficulty classes. We use SWE-bench Verified revision \code{c104f840cc67f8b6eec6f759ebc8b2693d585d4a}. Both arms use the same saved reports and the same Docker image for each pair. For each issue, we compare one baseline run with one TRACE run. We reused three such pairs from the pilot phase; for the remaining issues, we ran the baseline first and TRACE second. Configuration and report hashes are checked. The three pilot pairs recorded Docker image names rather than unique image fingerprints, limiting our ability to verify their original image contents. We have one selected generation per arm per instance, with earlier failed attempts preserved. No explicit generation seed was set. We did not repeat each method on each issue, so we cannot determine how much the results would change across repeated runs. Each assigned instance is the unit of analysis, not each generated test method. We record problems such as missing API credentials, network errors, and environment setup failures separately from rejected issue specifications. Fixing these setup problems does not permit changing the experimental instructions. For each issue and method, we use the generated patch identified in the saved run-selection record. We do not compare evaluation scores from earlier attempts and then choose the best-performing patch.

After test generation, we evaluate the saved patches using the official SWT-Bench harness at commit \code{eee38624e36a1b39976fffa2bddbc4c3c754d293}. The evaluator uses unit-test mode with standard patch handling and coverage collection enabled. It runs up to two evaluation jobs at once, with a 30-minute timeout for each execution. Both methods use identical evaluation settings. Dataset reference fields and evaluation outputs remain in an evaluation-only directory. All 26 assigned instances remain in the success-rate and patch-applicability denominators for each method. After infrastructure failures, recovery reuses unchanged environment-build recipes or predictions and preserves unsuccessful attempts. Candidates are not selected based on evaluation outcomes.

\section{Results}
\subsection{Success, Applicability, and Paired Outcomes}
Table~\ref{tab:results} reports official outcomes over all assignments. Baseline succeeds on 16/26 instances (61.54\%); TRACE succeeds on 17/26 (65.38\%). Both methods have 26/26 applicable patches (100\%): applicability means that the submitted patch applies cleanly to the repository. Table~\ref{tab:all} in Appendix~\ref{app:outcomes} lists every paired outcome and coverage value. Appendix~\ref{app:repro} provides reproducibility details, and Appendix~\ref{app:disclosure} records team contributions and AI-tool use. All 52 evaluation reports were produced---26 per method. No issues were removed from the totals because of technical problems. The evaluator reported no instance-level errors. Initial image timeouts and setup-check failures were resolved before final evaluation.

An evaluator exit code of zero indicates completed evaluation, not that every reproduction succeeded. The harness compares runs with the generated patch against a base comparison without that patch. It requires (i) additional fail-to-pass tests, (ii) preservation of the base comparison's fail-to-pass tests, and (iii) no additional tests that fail in both versions or change from passing to failing. Therefore, a patch can introduce a useful fail-to-pass test yet be unsuccessful because another assertion introduces a lasting regression. This stricter instance-level criterion explains some disagreements better than counting isolated failing assertions.

\begin{table}[t]
\caption{Official Evaluation and Selected Generation Records}
\label{tab:results}
\centering\small
\begin{tabular}{@{}p{0.48\columnwidth}rr@{}}
\toprule
Measure & Baseline & TRACE\\
\midrule
Official successes / assigned & 16/26 & 17/26\\
Success rate & 61.54\% & 65.38\%\\
Applicable / assigned & 26/26 & 26/26\\
Applicability rate & 100\% & 100\%\\
Mean coverage delta (cases with coverage) & 0.567364 & 0.602731\\
Coverage availability & 25/26 & 25/26\\
Mean coverage delta, successful cases & 0.691899 & 0.726739\\
Successful cases with coverage & 16 & 17\\
Mean generation time (seconds) & 93.45 & 127.09\\
Agent-only cost (USD) & 0.379072 & 0.383181\\
Evaluation time (seconds) & 3425.92 & 1549.86\\
\bottomrule
\end{tabular}
\end{table}

The paired outcomes are 15 instances where both succeed, one where only the baseline succeeds, two where only TRACE succeeds, and eight where neither succeeds (Table~\ref{tab:paired}). TRACE succeeds where the baseline does not on pytest 7432 and Django 15098; the baseline alone succeeds on Django 13344. The net difference is one instance, or 3.85 percentage points. The methods produced different success outcomes on only three issues: two favored TRACE and one favored the baseline. An exact two-sided McNemar test on the two TRACE-only and one baseline-only successes gives $p=1.0$. This small difference does not provide sufficient statistical evidence that either method performs better, nor does it establish that they perform equally.

\begin{table}[t]
\caption{Paired Official Reproduction Outcomes}
\label{tab:paired}
\centering\small
\begin{tabular}{lrr}
\toprule
 & TRACE succeeds & TRACE unsuccessful\\
\midrule
Baseline succeeds & 15 & 1\\
Baseline unsuccessful & 2 & 8\\
\bottomrule
\end{tabular}
\end{table}

TRACE accepts 22 issue specifications and falls back after three attempts on Django 15098, Django 16631, Sphinx 8551, and Django 13810. These fallbacks remain TRACE outcomes. Django 15098 is a TRACE-only success despite fallback, so that gain cannot be credited directly to the accepted issue specification. Django 11163 and Flask 5014 have byte-identical baseline/TRACE patches. TRACE accepts the issue specification for Django 13344 on its second extraction attempt, without fallback. Overall counts therefore describe the complete workflow, including fallback and sampling variation, rather than establishing whether specification extraction caused the difference.

\subsection{Coverage, Cost, and Runtime}
We report the official field \code{coverage_delta_pred}, the fraction of executable lines changed by the reference fix whose execution counts increase relative to the harness base comparison. This fraction is unitless. It is not whole-repository coverage. Arithmetic means exclude unavailable values and are not weighted by repository size or changed-line count. Overall means include all cases with available coverage, regardless of success. Success-only means use each method's successful cases with available coverage, so the two methods are averaged over different subsets.

Sphinx 8721 has N/A coverage in both arms because the harness finds no executable changed lines in its coverage inputs. N/A is excluded, not replaced with zero. The per-instance coverage values and aggregation counts are indexed in \code{evidence.json}; the metric implementation is in \code{src/grading.py} at the recorded SWT-Bench harness commit. A zero coverage delta means that no counted changed line has a positive execution-count increase relative to the base comparison; it does not mean that those lines were never executed. Success uses test outcomes rather than this coverage measure, so TRACE can succeed on Django 15098 and SymPy 24539 with zero delta. Baseline Django 15098 also has zero delta, but is unsuccessful. Identical Django 11163 patches receive different coverage deltas, showing that these measurements can vary even when submitted content matches.

Agent-reported costs for the selected runs total \$0.379072 for baseline and \$0.383181 for TRACE. TRACE additionally makes 69 preprocessing requests, recording 49,789 input and 173,755 output tokens. These costs exclude preprocessing and are not the total experiment cost. We report the recorded preprocessing token usage separately without assigning it a dollar amount. 

DeepSeek billing exports for October 18--November 16, 2026 show total charges of \$0.959976 (rounded) across the exported period. These charges are not limited to the selected 52 generation runs and cannot be separated reliably into baseline and TRACE totals from the daily billing records alone. They are reported separately, not added to the agent-reported costs, because adding them would double-count charges for the selected runs. 

Mean generation time recorded by the controller increases by 33.64 seconds, approximately 36\%, including TRACE preprocessing but excluding image preparation. Evaluation takes 57.10 minutes for baseline and 25.83 minutes for TRACE. TRACE runs second and reuses comparison caches, so evaluation timing is not evidence of workflow efficiency.

\section{Test Quality Analysis}\label{sec:test-quality}
Both the baseline and TRACE workflows instruct the agent to include meaningful assertions and inspect failures for syntax, import, setup, or issue-related causes. Execution records support manual review, but we did not implement an independent automatic checker for assertion quality or failure relevance.

We choose six distinct cases systematically: sort full instance IDs within each category, take the first two shared successes, first two shared failures, then first two disagreements, without duplication. This yields Astropy 14508, Django 11163, Django 12155, Django 13513, Django 13344, and Django 15098. Each category contains at least two cases, so no additional cases are needed. Appendix~\ref{app:cases} discusses the issue, generated check, both execution states, semantic relevance, and method implication for each.

The successful cases check clear expected behaviors: writing decimal numbers compactly in astronomy file headers without truncating comments, and returning an empty dictionary when no fields are requested. Both methods capture these behaviors; additional structure offers no demonstrated advantage there. The failures expose different weaknesses. For Django 12155, both methods produced assertions that exposed the indentation problem before the fix. However, the reference fix removed a helper function that the existing test module imported. This prevented the module from loading afterward, so neither method demonstrated a passing test on the fixed code. For Django 13513, both methods added tests checking that Django's debug report hides an earlier exception explicitly suppressed with \code{raise ... from None}. These tests passed before and after the fix because that behavior already worked. They did not exercise the scenario addressed by the official fix---an exception without a traceback---and therefore did not reproduce the bug. An accepted issue specification does not guarantee compatibility with the fixed code or the ability to distinguish buggy from fixed code.

The disagreement cases illustrate that an added persistent failure can make an instance unsuccessful even when new tests go from failing to passing. They also show that parsed test labels can differ from raw execution logs. Django 13344 shows that validating an issue specification does not guarantee a correct test setup: TRACE reproduced a similar symptom in a different middleware arrangement, but its test still failed after the fix. For Django 15098, both methods generated URL tests that failed before the fix and passed afterward. However, the baseline also added a language-lookup assertion that continued failing after the fix, causing its overall evaluation to fail. TRACE avoided that additional persistent failure and succeeded. Because TRACE used the original report after its issue specification failed validation, this result does not demonstrate a benefit from the extracted issue specification. Several official transition labels do not align cleanly with raw test-name lines; we report the official score and raw observations separately rather than silently correcting either.

\section{Threats to Validity}
Several factors limit our ability to attribute the observed improvement to TRACE. We selected one run per method for each issue, reused three pilot pairs, and ran the baseline before TRACE for the remaining issues. We did not measure variation across repeated runs or changes in the hosted model. The same model also extracted and reviewed the issue specification, so both steps could share the same misunderstanding.

Some recorded agent actions require additional review. Six selected runs were flagged for inspecting repository history: baseline on Django 13513, Django 15098, and Django 16631; TRACE on Seaborn 3187, Sphinx 8548, and Sphinx 8638. These are command-review flags, not proof that a reference fix was accessed or used. Other diagnostic actions included a temporary plugin that changed software behavior and stand-in code for missing dependencies. For Sphinx 8621, the baseline agent created a temporary pytest plugin that replaced the keyboard-formatting function in memory to try out its own proposed fix; passing tests under that plugin were not results from the official reference fix. For Sphinx 8638, the TRACE agent supplied a temporary replacement for the missing \code{roman} dependency, whose \code{toRoman()} function returned ordinary number strings instead of Roman numerals. This allowed diagnostic tests to run, but did not reproduce the real dependency's behavior. The submitted patches contain only test changes, but the patches alone do not establish that all actions during generation followed the experiment's restrictions.

The reported results depend on which tests the evaluator runs and how it reads their output. Some test names could not be matched between the before-fix and after-fix logs, and some recorded before/after transitions belonged to existing tests rather than newly generated ones. These issues can affect the reported success counts. Coverage measures execution of lines changed by the reference fix, not whether an assertion correctly represents the issue. We report the official results for all 26 issues and document discrepancies found in the raw logs. For three pilot pairs, we recorded Docker image names but not unique image fingerprints, so we cannot fully verify their original contents.

Our findings apply to a small set of assigned Python issues. Neither the Easy/Hard labels nor TRACE's one-instance advantage indicates how either method would perform on other projects. Provider billing totals are available for the exported period, but the records do not separate charges by workflow or selected run, so they cannot establish which method is more cost-effective.

A follow-up study should specify the number of repeated runs and how rule violations will be handled before reviewing any results. Human reviewers could also independently check whether each test captures the reported behavior and whether its failure comes from the bug rather than an environment problem. Measuring how often reviewers agree would provide an additional check beyond DeepSeek reviewing its own issue specification.

\section{Conclusion}
TRACE achieves 17/26 official successes versus 16/26 for the baseline, with 100\% patch applicability and increased mean generation time. The paired evidence supports a modest observed difference, not a reliable advantage for TRACE. Future work should repeat generation under recorded seeds where supported, compare extraction with and without review, check how raw test results are converted into parsed transitions, and add a shared quality checker for issue relevance and setup failures. Such a checker should initially assess both methods without using evaluator-only results to repair tests. Any later analysis of how the findings change under alternative assumptions should supplement the full assigned-set results rather than discard unfavorable cases.
